\section{Дослідження генерації асиметричних ключів}

\subsection{Генерація RSA-ключа довжиною 2048 біт}

\subsubsection{Програмна реалізація}

Для генерації RSA-ключа довжиною 2048 біт використано криптографічну бібліотеку \texttt{cryptography}. Час генерації вимірювався за допомогою функції \texttt{time.perf\_counter()}. Експеримент було проведено тричі.

\begin{lstlisting}[language=Python]
from cryptography.hazmat.primitives.asymmetric import rsa
import time

for i in range(1, 4):
    start = time.perf_counter()

    private_key = rsa.generate_private_key(
        public_exponent=65537,
        key_size=2048
    )

    end = time.perf_counter()
    print(f"Try (2048) {i}: {end - start:.4f} c")
\end{lstlisting}

\subsubsection{Результати генерації}

Під час трьох запусків час генерації RSA-ключа довжиною 2048 біт становив 0,0289 с, 0,0950 с і 0,0439 с відповідно. Середній час генерації дорівнював приблизно 0,0559 с.

\subsection{Генерація RSA-ключа довжиною 4096 біт}

\subsubsection{Програмна реалізація}

Для дослідження впливу довжини ключа на час генерації експеримент було повторено для RSA-ключа довжиною 4096 біт.

\begin{lstlisting}[language=Python]
from cryptography.hazmat.primitives.asymmetric import rsa
import time

for i in range(1, 4):
    start = time.perf_counter()

    private_key = rsa.generate_private_key(
        public_exponent=65537,
        key_size=4096
    )

    end = time.perf_counter()
    print(f"Try (4096) {i}: {end - start:.4f} c")
\end{lstlisting}

\subsubsection{Результати генерації}

Час генерації RSA-ключа довжиною 4096 біт становив 1,2870 с, 0,2607 с і 0,3652 с відповідно. Середній час генерації дорівнював приблизно 0,6376 с. Результати демонструють помітні відмінності між окремими запусками, що можуть бути пов'язані з особливостями пошуку простих чисел та навантаженням на систему.

\subsection{Порівняння часу генерації}

У таблиці наведено результати вимірювання часу генерації RSA-ключів різної довжини.

\begin{table}[H]
\centering
\caption{Результати генерації RSA-ключів}
\label{tab:rsa_generation}
\begin{tabular}{|c|c|c|}
\hline
\textbf{Довжина ключа, біт} &
\textbf{Номер запуску} &
\textbf{Час генерації, с} \\
\hline
2048 & 1 & 0,0289 \\
\hline
2048 & 2 & 0,0950 \\
\hline
2048 & 3 & 0,0439 \\
\hline
4096 & 1 & 1,2870 \\
\hline
4096 & 2 & 0,2607 \\
\hline
4096 & 3 & 0,3652 \\
\hline
\end{tabular}
\end{table}

Середній час генерації ключа довжиною 2048 біт становив приблизно 0,0559 с, а ключа довжиною 4096 біт — 0,6376 с. Таким чином, у проведеному експерименті генерація ключа довжиною 4096 біт у середньому тривала приблизно в 11,4 раза довше.

\subsection{Аналіз результатів}

Результати експерименту показують, що збільшення довжини RSA-ключа з 2048 до 4096 біт суттєво впливає на час генерації. Це пов'язано зі збільшенням розміру чисел, з якими виконуються арифметичні операції, та складністю пошуку відповідних простих чисел.

Час генерації може відрізнятися навіть для ключів однакової довжини. Під час кожного запуску використовуються нові випадкові значення, тому кількість спроб, необхідних для знаходження простих чисел, може бути різною. Крім того, на результати вимірювань впливають навантаження на процесор та інші процеси, що виконуються на комп'ютері.

Важливу роль у генерації RSA-ключів відіграє криптографічно стійкий генератор псевдовипадкових чисел (CSPRNG), який забезпечує непередбачуваність значень, що використовуються для формування кандидатів на прості числа. Ці кандидати проходять тести простоти, після чого знайдені прості числа використовуються для побудови ключової пари.

Збільшення довжини RSA-ключа також впливає на швидкість криптографічних операцій. Зокрема, операції розшифрування та створення цифрового підпису зазвичай потребують більше обчислень для довших ключів. Водночас збільшення довжини ключа підвищує стійкість RSA до атак, заснованих на факторизації модуля.

Самостійна реалізація генерації RSA-ключів у прикладному програмному забезпеченні зазвичай недоцільна, оскільки потребує коректної реалізації арифметики великих чисел, надійного джерела випадковості та перевірених методів тестування простоти. Помилки в цих компонентах можуть призвести до вразливості ключів. Тому доцільно використовувати перевірені криптографічні бібліотеки, які реалізують необхідні алгоритми та враховують вимоги безпеки.
